\subsection*{Formal Operational Inference Rules (Big-Step Semantics)}
We extend LO's big-step relation
$\langle e, \rho, H\rangle \Downarrow \langle v, H'\rangle$. A configuration
carries:
\begin{itemize}
  \item[--] the environment $\rho$, which maps variables to values;
  \item[--] the heap $H$, which maps objects to their class and fields;
  \item[--] fixed metadata $M$: each class's fields, methods and constructors,
    plus a table from class names to classes.
\end{itemize}

\textbf{Reflect Evaluation}
\[
\infrule{
  \langle \tm{expr}, \rho, H\rangle \Downarrow \langle \tm{object}, H_1\rangle \\
  \tm{object} \neq \tm{null} \\
  C = \tm{class}_{H_1}(\tm{object}) \\
  \tm{allocate\_info}(C, H_1) \Downarrow \langle \tm{info}, H_2\rangle
}{
  \langle \tm{reflect}(\tm{object}), \rho, H\rangle \Downarrow \langle \tm{info}, H_2\rangle
}
\]
{\small *\textbf{\code{allocate\_info(C, H)}} returns a fresh object q and the
heap H extended with $q \mapsto \tm{ReflectInfo}(C)$.}

\textbf{Field Read (\code{getfield})}
\[
\infrule{
  \langle \tm{expr}, \rho, H\rangle \Downarrow \langle \tm{object}, H_1\rangle \\
  \langle \tm{name\_expr}, \rho, H_1\rangle \Downarrow \langle \tm{name\_str}, H_2\rangle \\
  \tm{object} \neq \tm{null} \\
  M[\tm{class}_{H_2}(\tm{object})].\tm{field}(\tm{name\_str}) = (\tm{loc}, \tm{declared\_type}) \\
  \tm{declared\_type} \le T
}{
  \langle \tm{getfield}(T, \tm{expr}, \tm{name\_expr}), \rho, H\rangle \Downarrow
  \langle H_2[\tm{object}].\tm{fields}[\tm{loc}], H_2\rangle
}
\]

\textbf{Field Write (\code{setfield})}
\[
\infrule{
  \langle \tm{expr}, \rho, H\rangle \Downarrow \langle \tm{object}, H_1\rangle \\
  \langle \tm{name\_expr}, \rho, H_1\rangle \Downarrow \langle \tm{name\_str}, H_2\rangle \\
  \langle \tm{value\_expr}, \rho, H_2\rangle \Downarrow \langle \tm{new\_value}, H_3\rangle \\[1ex]
  \tm{object} \neq \tm{null} \\
  M[\tm{class}_{H_3}(\tm{object})].\tm{field}(\tm{name\_str}) = (\tm{loc}, \tm{declared\_type}) \\
  \tm{declared\_type} \le T \\
  \tm{new\_value} = \tm{null} \lor \tm{class}_{H_3}(\tm{new\_value}) \le \tm{declared\_type} \\
  H_4 = H_3[\tm{object}.\tm{fields}[\tm{loc}] \mapsto \tm{new\_value}]
}{
  \langle \tm{setfield}(T, \tm{expr}, \tm{name\_expr}, \tm{value\_expr});, \rho, H\rangle
  \Downarrow_{\tm{stmt}} \langle \rho, H_4\rangle
}
\]
{\small *The \textbf{\code{new\_value}} premise is needed because the static type
check only bounds the value by T; the runtime therefore checks the value's
actual class against the field's declared type.}

\textbf{Method Invocation Expression (\code{invoke})}
\[
\infrule{
  \langle \tm{expr}, \rho, H\rangle \Downarrow \langle \tm{object}, H_1\rangle \\
  \langle \tm{name\_expr}, \rho, H_1\rangle \Downarrow \langle \tm{name\_str}, H_2\rangle \\
  \langle \tm{arg\_expr}_i, \rho, H_{i+1}\rangle \Downarrow \langle \tm{arg\_value}_i, H_{i+2}\rangle
    \text{ for } i = 1, \ldots, k, \text{ in order} \\[1ex]
  \tm{object} \neq \tm{null} \\
  M[\tm{class}_{H_{k+2}}(\tm{object})].\tm{method}(\tm{name\_str}) =
    (\tm{method\_body}, [\tm{param}_1{:}\tm{param\_type}_1, \ldots, \tm{param}_m{:}\tm{param\_type}_m], \\
    \tm{declared\_type}) \\[1ex]
  \tm{declared\_type} = T, \quad T \neq \tm{void}, \quad k = m \\
  \tm{compatible}(\tm{arg\_type}_i, \tm{param\_type}_i) \text{ for every } i = 1, \ldots, k \\
  \tm{ordinary\_call}(\tm{method\_body}, \tm{object}, [\tm{arg\_value}_1, \ldots, \tm{arg\_value}_k], H_{k+2})
    \Downarrow \langle \tm{result}, H'\rangle
}{
  \langle \tm{invoke}(T, \tm{expr}, \tm{name\_expr}, \tm{arg\_expr}_1, \ldots, \tm{arg\_expr}_k), \rho, H\rangle
  \Downarrow \langle \tm{result}, H'\rangle
}
\]
{\small *\textbf{\code{compatible(arg\_type, param\_type)}} checks whether the
argument's recorded static type is assignment-compatible with the parameter
type.\\
*\textbf{\code{ordinary\_call(method\_body, object, args, H)}} executes the
selected method with \code{this} bound to object and parameters bound to args,
returning its result and updated heap.}

\textbf{Method Invocation Statement (\code{invoke} void)}
\[
\infrule{
  \langle \tm{expr}, \rho, H\rangle \Downarrow \langle \tm{object}, H_1\rangle \\
  \langle \tm{name\_expr}, \rho, H_1\rangle \Downarrow \langle \tm{name\_str}, H_2\rangle \\
  \langle \tm{arg\_expr}_i, \rho, H_{i+1}\rangle \Downarrow \langle \tm{arg\_value}_i, H_{i+2}\rangle
    \text{ for } i = 1, \ldots, k, \text{ in order} \\[1ex]
  \tm{object} \neq \tm{null} \\
  M[\tm{class}_{H_{k+2}}(\tm{object})].\tm{method}(\tm{name\_str}) =
    (\tm{method\_body}, [\tm{param}_1{:}\tm{param\_type}_1, \ldots, \tm{param}_m{:}\tm{param\_type}_m], \\
    \tm{declared\_type}) \\[1ex]
  \tm{declared\_type} = \tm{void}, \quad k = m \\
  \tm{compatible}(\tm{arg\_type}_i, \tm{param\_type}_i) \text{ for every } i = 1, \ldots, k \\
  \tm{ordinary\_call}(\tm{method\_body}, \tm{object}, [\tm{arg\_value}_1, \ldots, \tm{arg\_value}_k], H_{k+2})
    \Downarrow \langle \tm{unit}, H'\rangle
}{
  \langle \tm{invoke}(\tm{void}, \tm{expr}, \tm{name\_expr}, \tm{arg\_expr}_1, \ldots, \tm{arg\_expr}_k);, \rho, H\rangle
  \Downarrow_{\tm{stmt}} \langle \rho, H'\rangle
}
\]

\textbf{Dynamic Construction (\code{new\_instance})}
\[
\infrule{
  \langle \tm{name\_expr}, \rho, H\rangle \Downarrow \langle \tm{name\_str}, H_1\rangle \\
  \langle \tm{arg\_expr}_i, \rho, H_i\rangle \Downarrow \langle \tm{arg\_value}_i, H_{i+1}\rangle
    \text{ for } i = 1, \ldots, k, \text{ in order} \\[1ex]
  \tm{target\_class} = \tm{lookup\_user\_class}(\tm{name\_str}), \quad \tm{compatible}(\tm{target\_class}, T) \\
  \tm{constructor}(\tm{target\_class}, k) = \tm{selected\_ctor} \text{ with parameters }
    [\tm{param\_type}_1, \ldots, \tm{param\_type}_k] \\[1ex]
  \tm{compatible}(\tm{arg\_type}_i, \tm{param\_type}_i) \text{ for every } i = 1, \ldots, k \\[1ex]
  \tm{construct}(\tm{target\_class}, \tm{selected\_ctor}, [\tm{arg\_value}_1, \ldots, \tm{arg\_value}_k], H_{k+1})
    \Downarrow \langle \tm{object}, H'\rangle
}{
  \langle \tm{new\_instance}(T, \tm{name\_expr}, \tm{arg\_expr}_1, \ldots, \tm{arg\_expr}_k), \rho, H\rangle
  \Downarrow \langle \tm{object}, H'\rangle
}
\]
{\small *\textbf{\code{construct(target\_class, selected\_ctor, args, H)}}
allocates an instance, initializes all fields to LO defaults, then executes the
constructor and its delegation, returning the object and updated heap.}

\textbf{Dynamic Allocation (\code{allocate})}
\[
\infrule{
  \langle \tm{name\_expr}, \rho, H\rangle \Downarrow \langle \tm{name\_str}, H_1\rangle \\
  \tm{target\_class} = \tm{lookup\_user\_class}(\tm{name\_str}), \quad \tm{target\_class} \le T \\
  \tm{allocate\_default}(\tm{target\_class}, H_1) \Downarrow \langle \tm{object}, H_2\rangle
}{
  \langle \tm{allocate}(T, \tm{name\_expr}), \rho, H\rangle \Downarrow \langle \tm{object}, H_2\rangle
}
\]
{\small *\textbf{\code{allocate\_default(C, H)}} allocates an instance of C and
initializes all fields to LO defaults (the empty string for String fields),
without running a constructor, returning the object and updated heap.}


\subsection*{Metadata Layout}
Reflection metadata lives in read-only side tables
in \code{.rodata}, keyed by ClassDescriptor address, so ClassDescriptor and
the object header are unchanged. Per class, the tables record field names,
type tags and byte offsets (inherited first, in declaration order); method
names, return and parameter tags, and vtable slots (in vtable order); and
constructor arities and parameter tags, which only new\_instance uses. A
single class table maps each descriptor to its metadata and each class name
to its descriptor. The runtime searches it by linear scan, or sorts it by
descriptor address at initialization and uses binary search.

\textbf{Type tags.} Each type is one pointer-sized word: small integers for
void, int, bool, String and the null type, and the ClassDescriptor address
for a class (aligned, so never a small integer). An exact match is word
equality. A subtype check $A \le P$ holds if the words are equal, if $A$ is
the null type and $P$ is a class, or if both are classes and $P$ is on $A$'s
descriptor parent chain (the walk \code{lo\_instanceof} does). String is not
nullable in LO: null is assignable only to class types. Each reflective call
site passes a \code{.rodata} array of its arguments' static tags, with its
length, to the runtime helper; the argument values are passed separately,
already evaluated.
